\chapter{Templates and Metaprogramming}
\label{ch:c14-templates-metaprogramming}

\begin{quote}
\emph{C++14's template additions are small in number but large in leverage. Variable templates give constants a first-class parameterized form; \texttt{std::integer\_sequence} makes the compile-time ``indices trick'' a standard tool for unpacking tuples and parameter packs; the \texttt{\_t} trait aliases erase the \texttt{typename ...::type} noise that made C++11 metaprogramming unreadable; and \texttt{std::is\_final} closes a hole in the type-traits set.}
\end{quote}

Template metaprogramming in C++11 worked but read like assembly: every transformation trait needed \texttt{typename X<T>::type}, every compile-time constant was a \texttt{static constexpr} member buried in a struct, and unpacking a tuple required hand-rolled index machinery. C++14 polishes all three --- without adding a new metaprogramming paradigm --- making the same techniques dramatically more legible.

\section{Variable Templates}
\label{sec:c14-variable-templates}

Before C++14, a parameterized \emph{constant} had to be smuggled in as a static data member of a class template (\texttt{std::numeric\_limits<T>::max()}) or a \texttt{constexpr} function. \textbf{C++14 adds variable templates}: a variable can be templated directly, giving each type its own value with clean syntax.

\begin{lstlisting}[language=C++,caption={A variable template}]
template <typename T>
constexpr T pi = T(3.1415926535897932385);

// Each instantiation has the precision of its type:
float       pf = pi<float>;        // 3.14159f
double      pd = pi<double>;       // full double precision
long double pl = pi<long double>;  // extended precision
\end{lstlisting}

The same mechanism produces the now-ubiquitous \textbf{trait-value variables}: a boolean (or other) constant indexed by type, replacing \texttt{Trait<T>::value} with a plain \texttt{Trait\_v<T>}.

\begin{lstlisting}[language=C++,caption={Defining a trait-value variable template}]
template <typename T>
constexpr bool is_pointer_v = std::is_pointer<T>::value;

static_assert(is_pointer_v<int*>,  "");
static_assert(!is_pointer_v<int>,  "");
\end{lstlisting}

This is exactly the form the standard library adopted for its own \texttt{\_v} aliases (\texttt{std::is\_integral\_v}, etc.) --- \textbf{but those standard aliases are C++17}. In C++14 you have the \emph{language mechanism} (variable templates) and write the \texttt{\_v} helpers yourself, or keep using \texttt{Trait<T>::value}. The feature is the variable template; the standard \texttt{\_v} set is the C++17 application of it.

\section{\texttt{std::integer\_sequence} and the Indices Trick}
\label{sec:c14-integer-sequence}

\texttt{std::integer\_sequence<T, Is...>} is a compile-time list of integral constants encoded in a type. Its common specialization \texttt{std::index\_sequence<Is...>} (= \texttt{integer\_sequence<std::size\_t, Is...>}) carries a pack of indices \texttt{0, 1, ..., N-1}, generated by \texttt{std::make\_index\_sequence<N>}. Its purpose is to be a \textbf{bridge between the value world and the parameter-pack world}: it turns a count \texttt{N} into an expandable pack \texttt{0..N-1} you can use to subscript a tuple or array.

\begin{lstlisting}[language=C++,caption={The indices trick -- apply a function to an unpacked tuple}]
#include <tuple>
#include <utility>
#include <cstddef>

template <typename F, typename Tuple, std::size_t... Is>
auto apply_impl(F&& f, Tuple&& t, std::index_sequence<Is...>)
    -> decltype(std::forward<F>(f)(std::get<Is>(std::forward<Tuple>(t))...))
{
    // The pack Is... expands to 0, 1, ..., N-1, so this becomes
    //   f(get<0>(t), get<1>(t), ..., get<N-1>(t))
    return std::forward<F>(f)(std::get<Is>(std::forward<Tuple>(t))...);
}

template <typename F, typename Tuple>
auto apply(F&& f, Tuple&& t)
    -> decltype(apply_impl(std::forward<F>(f), std::forward<Tuple>(t),
                std::make_index_sequence<
                    std::tuple_size<std::decay_t<Tuple>>::value>{}))
{
    constexpr std::size_t N = std::tuple_size<std::decay_t<Tuple>>::value;
    return apply_impl(std::forward<F>(f), std::forward<Tuple>(t),
                      std::make_index_sequence<N>{});
}
\end{lstlisting}

\begin{lstlisting}[language=C++,caption={Calling it}]
int add3(int a, int b, int c) { return a + b + c; }
auto args = std::make_tuple(1, 2, 3);
int result = apply(add3, args);   // -> add3(1, 2, 3) == 6
\end{lstlisting}

The pattern is always the same two-function shape: a public function computes the size and manufactures an \texttt{index\_sequence}, and a private \texttt{\_impl} overload accepts the sequence purely to \textbf{name the pack \texttt{Is...}} that drives the expansion. Beyond tuple unpacking, \texttt{make\_index\_sequence} powers compile-time loops over packs and array initialization from a generator. (\texttt{std::index\_sequence\_for<Ts...>} is shorthand for \texttt{make\_index\_sequence<sizeof...(Ts)>} when you already have a type pack.)

\begin{quote}
\textbf{Note:} \texttt{std::apply} itself became a \emph{standard} function in \textbf{C++17}; in C++14 you write the \texttt{apply\_impl}/\texttt{apply} pair above by hand. The reusable primitive C++14 gives you is \texttt{integer\_sequence}.
\end{quote}

\section{Transformation Trait Aliases (\texttt{\_t})}
\label{sec:c14-trait-aliases}

C++11 transformation traits returned their result as a nested \texttt{::type}, forcing a \texttt{typename} disambiguator at every use. A SFINAE constraint or a type transformation therefore drowned in \texttt{typename ...::type} noise. \textbf{C++14 adds \texttt{\_t} alias templates} for every standard transformation trait, collapsing \texttt{typename Trait<T>::type} to \texttt{Trait\_t<T>}.

\begin{lstlisting}[language=C++,caption={\_t aliases vs the C++11 spelling}]
// C++11:
typename std::enable_if<std::is_integral<T>::value, T>::type f(T);
typename std::remove_reference<T>::type        x;
typename std::decay<T>::type                   y;

// C++14:
std::enable_if_t<std::is_integral<T>::value, T> f(T);
std::remove_reference_t<T>                       x;
std::decay_t<T>                                  y;
\end{lstlisting}

The aliases are mechanical --- \texttt{template <class T> using remove\_reference\_t = typename remove\_reference<T>::type;} --- but the readability gain is decisive in the place they matter most, SFINAE-constrained signatures:

\begin{lstlisting}[language=C++,caption={SFINAE made readable with enable\_if\_t}]
template <typename T,
          typename = std::enable_if_t<std::is_floating_point<T>::value>>
T reciprocal(T x) { return T(1) / x; }
\end{lstlisting}

Every standard trait with a \texttt{::type} result has a corresponding \texttt{\_t} alias (\texttt{decay\_t}, \texttt{remove\_reference\_t}, \texttt{remove\_cv\_t}, \texttt{add\_pointer\_t}, \texttt{conditional\_t}, \texttt{common\_type\_t}, \texttt{underlying\_type\_t}, \texttt{result\_of\_t}, and the rest). Prefer them universally --- they are pure noise reduction with no semantic cost. (The parallel \texttt{\_v} value aliases for the \emph{predicate} traits, by contrast, are C++17; see \S\ref{sec:c14-variable-templates}.)

\section{\texttt{std::is\_final}}
\label{sec:c14-is-final}

C++14 added the type trait \textbf{\texttt{std::is\_final<T>}}, which reports whether a class type was declared \texttt{final}. It fills a gap in the C++11 traits set and lets generic code detect a sealed class --- for instance, to decide whether deriving from it (a common empty-base-optimization technique) is even legal.

\begin{lstlisting}[language=C++,caption={Detecting a final class}]
#include <type_traits>

struct Sealed final {};
struct Open {};

static_assert(std::is_final<Sealed>::value,  "Sealed is final");
static_assert(!std::is_final<Open>::value,   "Open is not final");
\end{lstlisting}

A concrete use is guarding the \emph{derive-to-compress} idiom (empty base optimization): a wrapper that inherits from its policy type to elide storage for a stateless policy cannot inherit from a \texttt{final} policy, so the trait selects between the inheriting and the member-holding implementation.

\begin{lstlisting}[language=C++,caption={Choosing storage strategy based on is\_final}]
template <typename Policy,
          bool CanDerive = !std::is_final<Policy>::value &&
                            std::is_empty<Policy>::value>
class Holder;   // primary template

// CanDerive == true  -> inherit from Policy (empty base optimization)
// CanDerive == false -> hold Policy as a member
\end{lstlisting}

As with the other predicate traits, the value is read through \texttt{::value} in C++14; the \texttt{std::is\_final\_v} shorthand is C++17.

\section{Professional Insights}
\label{sec:c14-templates-insights}

\textbf{Use variable templates for parameterized constants and your own trait shorthands.} A \texttt{template <class T> constexpr T epsilon = ...;} reads better than a function call or a buried static member, and defining \texttt{Trait\_v<T>} helpers locally gives you the clean \texttt{\_v} ergonomics in a C++14 codebase years before C++17's standard aliases. Keep them \texttt{constexpr} so they fold to immediates.

\textbf{\texttt{integer\_sequence} is the standard answer to ``I have a count, I need a pack.''} Whenever you must expand over \texttt{0..N-1} --- unpacking a tuple into a call, initializing an array from a generator, looping over a parameter pack at compile time --- reach for \texttt{make\_index\_sequence} and the two-function \texttt{\_impl} pattern rather than re-deriving bespoke index machinery.

\textbf{Adopt the \texttt{\_t} aliases everywhere --- there is no downside.} They are pure syntactic relief, and in SFINAE-heavy template code the difference between \texttt{typename std::enable\_if<C, T>::type} and \texttt{std::enable\_if\_t<C, T>} is the difference between a signature you can read and one you can't. Make their use a style rule.

\textbf{Mind the version line on \texttt{\_v} and \texttt{std::apply}.} Variable templates, \texttt{integer\_sequence}, the \texttt{\_t} aliases, and \texttt{is\_final} are all C++14; the \emph{standard} \texttt{\_v} predicate aliases (\texttt{is\_integral\_v}) and \texttt{std::apply} are C++17. In a strict C++14 target, define the \texttt{\_v} helpers and the \texttt{apply} wrapper yourself --- the underlying mechanisms are already present.
